% paper/figs/fig3_caption.tex -- GENERATED by the fig3 driver.
% Do not edit: rebuild with paper/build.sh.
\newcommand{\figthreecaption}{%
The double-helicity-flip structure function extracted from the amplitudes of Fig.~\ref{fig:inclusive}, at $e(10\,\mathrm{GeV}) \times {}^{6}\mathrm{Li}\,(99.5\,\mathrm{GeV/u})$ with $P_{zz} = 0.60$: $x\Delta(x, Q^{2})$ per nucleon in the $Q^{2} = 1.14$\,GeV$^{2}$ (a) and $Q^{2} = 14.3$\,GeV$^{2}$ (b) slices, with independent one-year (open, 10\,fb$^{-1}$ per nucleon) and ten-year (filled, 100\,fb$^{-1}$ per nucleon) pseudo-measurements against the injected moment-constrained $\Delta$ and the no-$F_{1}$ variant. Each point is the fitted amplitude of a pair of merged $x$ bins converted to a point value by the model bin-centring factor $K = \Delta_{\rm model}(x_{c}, Q^{2}_{c})/A_{\rm model}$, so that $\delta\Delta = \delta\hat A\,|K|$. The best bin of each slice reaches $\Delta = -0.135 \pm 0.003$ at $x = 0.02$ and $-0.0047 \pm 0.0004$ at $x = 0.316$ in one year, 2.5\% and 9.5\% relative. The two curves are separated bin by bin at low $Q^{2}$ and cross at large $x$ in the high-$Q^{2}$ slice, which is what makes the $x$ dependence, rather than the size of the amplitude, the discriminating measurement. Statistical errors only; the third sweet-spot slice, $Q^{2} = 3.13$\,GeV$^{2}$, is measured with the same draw and is not drawn here.%
}
